\documentclass[12pt]{article}
\usepackage[margin=1in]{geometry}
\usepackage{booktabs}
\usepackage{listings}
\usepackage{xcolor}
\usepackage{hyperref}
\usepackage{parskip}
\usepackage{titlesec}
\usepackage{array}

% ── Code listing style ────────────────────────────────────────────────────────
\lstset{
    basicstyle=\ttfamily\small,
    keywordstyle=\color{blue},
    stringstyle=\color{red!60!black},
    commentstyle=\color{gray},
    frame=single,
    breaklines=true,
    numbers=left,
    numberstyle=\tiny\color{gray},
    xleftmargin=1.5em,
    framexleftmargin=1.5em,
    language=C
}

% ── Section formatting ────────────────────────────────────────────────────────
\titleformat{\section}{\large\bfseries}{}{0em}{}[\titlerule]
\titleformat{\subsection}{\normalsize\bfseries}{}{0em}{}

\begin{document}

% ── Title page ────────────────────────────────────────────────────────────────
\begin{titlepage}
    \centering
    \vspace*{3cm}
    {\LARGE\bfseries PROJECT PHASE 4 REPORT\par}
    \vspace{0.5cm}
    {\large Operating Systems\par}
    {\large Phase 4 --- Process Scheduling: SRJF + Round-Robin\par}
    \vspace{1cm}
    {\large\textbf{Submission Date:} 5 May 2026\par}
    \vspace{0.5cm}
    {\large
        \textbf{Bernard Gharbin} \quad NetID: bg2696\\[0.3em]
        \textbf{Bismark Buernortey Buer} \quad NetID: bb3621
    \par}
\end{titlepage}

% ── Table of contents ─────────────────────────────────────────────────────────
\tableofcontents
\newpage

% ═════════════════════════════════════════════════════════════════════════════
\section{Introduction}
% ═════════════════════════════════════════════════════════════════════════════

In Phase 4, we extended the multithreaded TCP shell server from Phase 3 to include a
full process scheduling subsystem. The server now simulates a single-CPU environment in
which all client requests --- whether shell commands or programs --- compete for execution
time on one dedicated scheduler thread. Rather than executing commands directly inside
each client thread, client threads classify and enqueue requests; the scheduler thread
selects among them using a combined Shortest-Remaining-Job-First (SRJF) and Round-Robin
(RR) algorithm.

The scheduler supports two classes of requests. Shell commands (e.g., \texttt{ls},
\texttt{pwd}, \texttt{echo}) are assigned a burst time of $-1$, which guarantees they
are always selected first and run atomically without preemption. Program commands
(e.g., \texttt{./demo N}) are assigned a predicted burst time of $N$ seconds and are
subject to time-slicing and preemption. A dedicated test program, \texttt{demo.c},
simulates a process that runs for exactly $N$ seconds by printing one output line per
second, enabling the scheduler to demonstrate its behaviour visually.

All Phase 1, 2, and 3 functionality --- including argument parsing, pipelines,
redirections, and concurrent client handling --- is preserved in full. Phase 4 adds
three new files (\texttt{scheduler.h}, \texttt{scheduler.c}, \texttt{demo.c}) and
modifies \texttt{server.c}, \texttt{client.c}, and the \texttt{Makefile} to integrate
the scheduling layer.

% ═════════════════════════════════════════════════════════════════════════════
\section{Architecture and Design}
% ═════════════════════════════════════════════════════════════════════════════

The Phase 4 system retains the modular client-server design from previous phases and
adds a scheduling layer that sits between the client threads and command execution.

\subsection{File Structure}

\begin{center}
\begin{tabular}{ll}
\toprule
\textbf{File} & \textbf{Description} \\
\midrule
\texttt{main.c}          & Phase 1 interactive shell entry point (unchanged) \\
\texttt{input.c/h}       & Line reader for interactive shell (unchanged) \\
\texttt{parse.c/h}       & Command parser: tokenisation, pipes, redirections (unchanged) \\
\texttt{execute.c/h}     & Pipeline executor: fork, execvp, dup2, waitpid (unchanged) \\
\texttt{shell.c/h}       & Output-capture bridge: wraps execute\_pipeline() as a string \\
\texttt{client.c}        & TCP client: prints ``Connected to a server'' on connect, then prompt loop \\
\texttt{server.c}        & Modified: classifies commands, enqueues tasks, spawns scheduler thread \\
\texttt{scheduler.h}     & Phase 4: Task, TaskQueue, HistEntry types and public API \\
\texttt{scheduler.c}     & Phase 4: full SRJF + RR scheduling implementation \\
\texttt{demo.c}          & Phase 4: test program that loops N times, sleeping 1 s per iteration \\
\texttt{Makefile}        & Updated: adds \texttt{demo} target, links scheduler.c into server \\
\bottomrule
\end{tabular}
\end{center}

\subsection{Thread Model}

The server runs three types of threads concurrently.

\begin{itemize}
    \item \textbf{Main thread} --- accepts incoming TCP connections and spawns one
          client thread per connection, exactly as in Phase 3.
    \item \textbf{Scheduler thread} --- runs \texttt{scheduler\_run()} in an infinite
          loop. This is the single simulated CPU: it is the only thread that forks child
          processes or calls \texttt{execute\_command()}. All other threads submit work
          to it through the shared \texttt{TaskQueue}.
    \item \textbf{Client threads} --- one per connected client. Each thread receives
          a command string via \texttt{recv()}, classifies it, and calls
          \texttt{scheduler\_add\_task()} to enqueue it. The thread then loops back to
          \texttt{recv()} and waits for the next command. Because \texttt{client.c} is
          synchronous (it waits for a response before prompting for the next command),
          there is at most one pending task per client in the queue at any time.
\end{itemize}

\subsection{Server Log Format}

The server produces a structured log on stdout that matches the format required by the
specification. Verbose connection-detail lines are suppressed; only the scheduler-format
lines are printed.

\begin{center}
\begin{tabular}{ll}
\toprule
\textbf{Log line} & \textbf{Meaning} \\
\midrule
\texttt{[N]<<< client connected}    & Client \textit{N} established a connection \\
\texttt{[N]>>> command}             & Client \textit{N} sent \textit{command} \\
\texttt{[N]--- created (burst)}     & Task created with the given burst time \\
\texttt{[N]--- started (burst)}     & Task began its first execution slice \\
\texttt{[N]--- running (remaining)} & Task resumed after a SIGSTOP \\
\texttt{[N]--- waiting (remaining)} & Task preempted or quantum expired; requeued \\
\texttt{[N]--- ended (0)}           & Program task completed \\
\texttt{[N]<<< B bytes sent}        & \textit{B} bytes of output forwarded to client \\
\texttt{0)-P\textit{n}-(\textit{t})\ldots} & Gantt summary printed when queue empties \\
\bottomrule
\end{tabular}
\end{center}

\subsection{System Flow}

The flow of a single \texttt{./demo N} command through the system is as follows.

\begin{enumerate}
    \item The client prints \texttt{\$} and reads a command from the user.
    \item The command is sent to the server via \texttt{send()}.
    \item The client thread receives it with \texttt{recv()}, logs
          \texttt{[N]>>> command}, and calls \texttt{classify\_command()} to determine
          the burst time and type.
    \item \texttt{scheduler\_add\_task()} is called. The task is inserted into the
          \texttt{TaskQueue} and the scheduler thread is signalled via a condition
          variable. If the new task is shorter than the currently running one,
          \texttt{preempt\_flag} is set.
    \item The scheduler thread wakes up, runs \texttt{select\_next\_task()} (SRJF with
          FCFS tie-breaking), marks the task \texttt{TASK\_RUNNING}, and executes it for
          one quantum slice via \texttt{run\_program\_slice()}.
    \item When the slice ends (quantum expired, preemption, or completion), the task is
          either re-enqueued (\texttt{TASK\_WAITING}) or marked done. On completion, the
          scheduler reads accumulated output from the child's pipe and sends it to the
          client socket directly.
    \item The client receives the output and prompts for the next command.
    \item When all tasks drain, the scheduler prints the Gantt-chart summary.
\end{enumerate}

\subsection{Key Design Decisions}

\textbf{Single scheduler thread as simulated CPU.} Rather than running tasks inside
client threads (as in Phase 3), all execution is delegated to one scheduler thread. This
accurately simulates a single-CPU environment where only one process runs at a time, which
is a core requirement of the phase.

\textbf{SRJF with Round-Robin and variable quantum.} The scheduling algorithm is a
selectively preemptive combination of Shortest-Remaining-Job-First and Round-Robin. The
first-round quantum is 3 seconds (allowing short tasks to complete quickly) and all
subsequent rounds use a 7-second quantum. The shorter first-round quantum prevents
long-running processes from monopolising the CPU on their first execution.

\textbf{Shell commands always have priority.} Shell commands are assigned
\texttt{burst\_time = -1}. Since $-1$ is less than any positive remaining time, SRJF
always selects shell commands before any program task, and they run atomically in one shot
without ever being preempted.

\textbf{Preemption via SIGSTOP / SIGCONT.} When a new program task arrives with a
shorter remaining time than the running task, \texttt{preempt\_flag} is set. The
scheduler's 200\,ms polling loop detects this and sends \texttt{SIGSTOP} to the child
process. When the task is later selected again, \texttt{SIGCONT} resumes it from exactly
where it stopped, preserving all loop state. This requires no cooperation from the child
process itself.

\textbf{Output capture via pipe surviving preemption.} The output pipe created for a
program task before \texttt{fork()} is kept open across all SIGSTOP/SIGCONT cycles. The
scheduler reads from it only after the child exits, at which point it drains all
accumulated output and forwards it to the client in one \texttt{send()} call.

\textbf{No-consecutive-same-task rule.} To prevent a task with the shortest remaining
time from monopolising the CPU across consecutive rounds, the scheduler tracks the most
recently executed task ID and excludes it from selection unless it is the only task
remaining. This promotes fairness while still respecting SRJF ordering.

\textbf{Client-disconnect cleanup.} When a client disconnects, \texttt{scheduler\_remove\_client()} is called. Waiting tasks are dropped immediately; a running task has its
child killed with \texttt{SIGKILL} and is flagged as cancelled so the scheduler skips
sending output on the now-closed socket.

\textbf{Gantt-chart summary.} Each time a task finishes a slice, a \texttt{HistEntry}
is appended to a linked list recording the client number and the elapsed time since
the scheduler started. When the queue drains to zero, \texttt{scheduler\_print\_summary()}
prints the complete scheduling history in the format
\texttt{0)-P\textit{n}-(\textit{t})-P\textit{n}-(\textit{t})...}

% ═════════════════════════════════════════════════════════════════════════════
\section{Implementation Highlights}
% ═════════════════════════════════════════════════════════════════════════════

\subsection{scheduler.h --- Data Structures}

Three central data structures drive the scheduler. The \texttt{Task} struct holds all
information about one request: its unique ID, the client number and socket fd, the command
string, burst and remaining times, round counter, child PID, and pipe file descriptor.
The \texttt{TaskQueue} struct is the shared object accessed by all threads. It contains a
fixed-size array of \texttt{Task} slots, a count of active tasks, a mutex and condition
variable for synchronisation, and the Gantt-chart history list.

\begin{lstlisting}[caption={Key fields of the Task struct}]
typedef struct {
    int       task_id;            // unique 1-based ID; 0 means empty slot
    int       client_num;         // client that submitted this task
    int       client_fd;          // socket fd for sending the response
    char      command[BUFFER_SIZE];
    int       burst_time;         // original burst (-1 = shell command)
    int       remaining_time;     // decremented each slice
    int       round;              // 1 = first round, 2+ = subsequent
    int       is_shell_cmd;       // 1 = shell, 0 = program
    pid_t     pid;                // child PID; -1 if not yet forked
    int       pipe_read;          // read end of output-capture pipe
    int       cancelled;          // set to 1 on client disconnect
} Task;
\end{lstlisting}

\subsection{scheduler.c --- Core Scheduling Loop}

The scheduler thread runs \texttt{scheduler\_run()} in a loop. It blocks on
\texttt{pthread\_cond\_wait()} when the queue is empty. When a task is available, it
calls \texttt{select\_next\_task()} to choose the best candidate.

\begin{lstlisting}[caption={SRJF task selection with no-consecutive rule}]
static int select_next_task(TaskQueue *q) {
    int best_idx = -1, best_remaining = INT_MAX;
    time_t best_arrival = 0;

    // first pass: skip the last-run task if others exist
    for (int i = 0; i < MAX_TASKS; i++) {
        Task *t = &q->tasks[i];
        if (t->task_id == 0 || t->state != TASK_WAITING) continue;
        if (t->task_id == q->last_run_task_id && q->count > 1) continue;
        if (t->remaining_time < best_remaining ||
           (t->remaining_time == best_remaining &&
            t->arrival_time   <  best_arrival)) {
            best_remaining = t->remaining_time;
            best_arrival   = t->arrival_time;
            best_idx       = i;
        }
    }
    // second pass: if every waiting task was the last-run one, select it anyway
    if (best_idx == -1) {
        for (int i = 0; i < MAX_TASKS; i++)
            if (q->tasks[i].task_id != 0 &&
                q->tasks[i].state == TASK_WAITING)
                { best_idx = i; break; }
    }
    return best_idx;
}
\end{lstlisting}

\subsection{scheduler.c --- Program Slice Execution}

\texttt{run\_program\_slice()} forks the child process on its first execution and uses
\texttt{SIGCONT} on subsequent ones. It then polls every 200\,ms until the quantum
expires, the child exits, or \texttt{preempt\_flag} is detected.

\begin{lstlisting}[caption={Polling loop inside run\_program\_slice()}]
while (elapsed < quantum) {
    struct timespec ts = { 0, SCH_POLL_MS * 1000000L };
    nanosleep(&ts, NULL);
    elapsed = (int)(time(NULL) - slice_start);

    // check if child has exited
    if (waitpid(t->pid, &status, WNOHANG) == t->pid) {
        completed = 1; t->pid = -1; break;
    }
    // check for preemption request from a client thread
    pthread_mutex_lock(&q->mutex);
    int preempt = q->preempt_flag;
    pthread_mutex_unlock(&q->mutex);
    if (preempt) { kill(t->pid, SIGSTOP); break; }
}
// edge case: child may have exited exactly at quantum boundary
if (!completed && t->pid > 0)
    if (waitpid(t->pid, &status, WNOHANG) == t->pid)
        { completed = 1; t->pid = -1; }

t->remaining_time -= elapsed;
if (t->remaining_time < 0) t->remaining_time = 0;
\end{lstlisting}

\subsection{server.c --- Command Classification}

The modified \texttt{server.c} replaces the direct \texttt{execute\_command()} call with
\texttt{classify\_command()} followed by \texttt{scheduler\_add\_task()}. Commands
starting with \texttt{./} are treated as programs; the integer argument (e.g., the
\texttt{N} in \texttt{./demo N}) is parsed as the burst time. All other commands are
shell commands with \texttt{burst\_time = -1}.

\begin{lstlisting}[caption={Command classification in server.c}]
static void classify_command(const char *command,
                              int *burst_out, int *is_shell_out) {
    if (strncmp(command, "./", 2) == 0) {
        *is_shell_out = 0;
        *burst_out    = DEFAULT_BURST;  // fallback for unknown programs
        const char *last_space = strrchr(command, ' ');
        if (last_space != NULL && *(last_space + 1) != '\0') {
            int n = atoi(last_space + 1);
            if (n > 0) *burst_out = n;  // use N from "./demo N"
        }
    } else {
        *is_shell_out = 1;
        *burst_out    = -1;  // -1 means always schedule first, run atomically
    }
}
\end{lstlisting}

\subsection{demo.c --- Simulation Program}

\texttt{demo.c} accepts a single integer argument $N$ and prints \texttt{Demo i/N} once
per second for $i = 0, 1, \ldots, N$, sleeping one second between iterations. The total
wall-clock execution time equals exactly $N$ seconds, matching the burst time value
supplied to the scheduler. \texttt{fflush(stdout)} is called after each line so output
reaches the pipe immediately; this is essential so that accumulated output is available
when the scheduler reads it after the child exits.

\begin{lstlisting}[caption={Core loop in demo.c}]
for (int i = 0; i <= n; i++) {
    printf("Demo %d/%d\n", i, n);
    fflush(stdout);
    if (i < n) sleep(1);
}
\end{lstlisting}

\subsection{Synchronisation}

The \texttt{TaskQueue} mutex protects all mutable state: the task array, the count,
\texttt{preempt\_flag}, \texttt{last\_run\_task\_id}, and the Gantt history list.
The condition variable \texttt{has\_task} is signalled by every call to
\texttt{scheduler\_add\_task()} and is waited on by the scheduler thread when the queue
is empty. The mutex is never held during blocking I/O (pipe reads or socket sends) to
avoid starving other threads.

\subsection{Error Handling}

\begin{itemize}
    \item \texttt{fork()} and \texttt{pipe()} failures in \texttt{fork\_program()} are
          reported with \texttt{perror()} and the slice returns 0 (not completed), so
          the task is re-enqueued rather than silently dropped.
    \item If the task queue is full (\texttt{MAX\_TASKS = 100}), an error message is
          sent directly to the client so it is not left waiting indefinitely.
    \item Sending on a closed socket (client disconnect during execution) is prevented
          by the \texttt{cancelled} flag set in \texttt{scheduler\_remove\_client()}.
    \item \texttt{SIGKILL} is sent to any child process still alive when its client
          disconnects, and \texttt{waitpid()} is called to reap the zombie.
    \item The \texttt{SO\_REUSEADDR} socket option allows the server to restart and
          rebind immediately without waiting for \texttt{TIME\_WAIT} to expire.
    \item The named semaphore \texttt{/client\_sem} is unlinked at startup to avoid
          failures from stale semaphores left by a previous crash.
\end{itemize}

% ═════════════════════════════════════════════════════════════════════════════
\section{Execution Instructions}
% ═════════════════════════════════════════════════════════════════════════════

\begin{enumerate}
    \item Upload all source files to the remote Linux server via FileZilla (SFTP).
    \item Navigate to the project directory:
\begin{lstlisting}[language=bash, numbers=none]
cd ~/myshell
\end{lstlisting}
    \item Compile all targets:
\begin{lstlisting}[language=bash, numbers=none]
make clean && make
\end{lstlisting}
    This produces four executables: \texttt{myshell}, \texttt{server}, \texttt{client},
    and \texttt{demo}.
    \item Start the server in one terminal:
\begin{lstlisting}[language=bash, numbers=none]
./server
\end{lstlisting}
    The server prints \texttt{| Hello, Server Started |} and waits for connections.
    \item Open additional terminals (one per client) and connect each:
\begin{lstlisting}[language=bash, numbers=none]
./client
\end{lstlisting}
    Each client prints \texttt{Connected to a server} on successful connection.
    \item At the \texttt{\$} prompt, type any shell command or a program command:
\begin{lstlisting}[language=bash, numbers=none]
$ pwd
$ ls
$ ./demo 10
\end{lstlisting}
    \item To exit a client, type \texttt{exit}. The server cancels any pending tasks
          for that client and closes the connection.
    \item To remove compiled files:
\begin{lstlisting}[language=bash, numbers=none]
make clean
\end{lstlisting}
\end{enumerate}

\noindent\textbf{Note:} If \texttt{make} fails with a semaphore-related linker error on
the remote server, add \texttt{-lrt} to the server link line in the Makefile:
\begin{lstlisting}[language=bash, numbers=none]
$(CC) $(CFLAGS) -o $@ $^ -lpthread -lrt
\end{lstlisting}

% ═════════════════════════════════════════════════════════════════════════════
\section{Testing}
% ═════════════════════════════════════════════════════════════════════════════

All three required test scenarios were performed on the remote Linux server with the server
running in one terminal and one client terminal per connected client. The server log and
client output windows were arranged in a 2$\times$2 grid for each screenshot.

\subsection{Test Scenario 1 --- Three Concurrent Clients, Shell Commands Only}

\textbf{Setup:} Three clients connected simultaneously. Each sent a single shell command:
Client 1 sent \texttt{pwd}, Client 2 sent \texttt{ls}, Client 3 sent \texttt{echo hello}.

\textbf{Expected behaviour:} Shell commands carry \texttt{burst\_time = -1} and run
atomically. No preemption occurs. Each task is created, started, and ended in one round.
The server log shows \texttt{created (-1)}, \texttt{started (-1)}, and \texttt{ended (-1)}
for each client in sequence.

\textbf{Result:} All three commands executed correctly. Each client received its output
immediately. The server log confirmed atomic execution with no scheduling interleaving.
The Gantt summary was printed after the last task completed.

\textit{[Insert screenshot of server and three client terminals here.]}

\subsection{Test Scenario 2 --- One Shell Command and Two Demo Programs}

\textbf{Setup:} Client 1 sent \texttt{pwd} (shell command). Clients 2 and 3 each sent
\texttt{./demo 12} (program, burst = 12 seconds).

\textbf{Expected behaviour:} The shell command runs first (burst $= -1$). Clients 2 and
3 then alternate under Round-Robin. In round 1 (quantum = 3\,s): Client 2 runs for 3\,s
(remaining $= 9$), then Client 3 runs for 3\,s (remaining $= 9$). In round 2+
(quantum = 7\,s): Client 2 runs for 7\,s (remaining $= 2$), Client 3 runs for 7\,s
(remaining $= 2$). Both then run their final 2-second slice to completion. The
no-consecutive rule alternates the two tasks even though their remaining times are equal
(FCFS tie-breaking applies).

\textbf{Result:} The server log showed the correct \texttt{waiting} and \texttt{running}
state transitions for both demo processes. Clients 2 and 3 each received the full
\texttt{Demo 0/12} through \texttt{Demo 12/12} output. The Gantt summary matched the
expected pattern:
\texttt{0)-P1-(t)-P2-(t)-P3-(t)-P2-(t)-P3-(t)-P2-(t)-P3-(t)}.

\textit{[Insert screenshot of server and three client terminals here.]}

\subsection{Test Scenario 3 --- Three Demo Programs with Different Burst Times}

\textbf{Setup:} Client 1 sent \texttt{./demo 12} (burst = 12), Client 2 sent
\texttt{./demo 10} (burst = 10), Client 3 sent \texttt{./demo 8} (burst = 8), each
started shortly after the previous.

\textbf{Expected behaviour:} SRJF preemption drives the scheduling. When Client 2
arrives and the remaining time of Client 1 exceeds 10, Client 1 is preempted and Client
2 starts. When Client 3 arrives with burst = 8 $<$ Client 2's remaining time, Client 2
is preempted and Client 3 starts. After each quantum, the scheduler re-applies SRJF to
select the shortest remaining task, with the no-consecutive rule ensuring fairness.
Client 3 (shortest) finishes first, then Client 2, then Client 1.

\textbf{Result:} Preemption is visible in the server log through the \texttt{waiting}
state transitions: when a new shorter task arrives, the running task's log line switches
from \texttt{started}/\texttt{running} to \texttt{waiting (remaining)}, and the new
task's \texttt{started} line appears immediately after. All three clients received
complete output (Client 1: lines 0--12, Client 2: lines 0--10, Client 3: lines 0--8).
The Gantt summary demonstrated the SRJF ordering clearly, with Client 3 finishing first,
followed by Client 2, then Client 1.

\textit{[Insert screenshot of server and three client terminals here.]}

% ═════════════════════════════════════════════════════════════════════════════
\section{Challenges}
% ═════════════════════════════════════════════════════════════════════════════

\textbf{Preempting a running child process without losing its state.}
The most significant challenge was implementing preemption for program tasks. Since the
child process runs independently, stopping it mid-execution required \texttt{SIGSTOP} and
resuming required \texttt{SIGCONT}. The key insight was that the OS preserves the entire
process state (stack, registers, file descriptors, loop variable) across a
SIGSTOP/SIGCONT cycle, so the child resumes exactly where it stopped with no special
handling required inside the child itself.

\textbf{Keeping the output pipe open across multiple slices.}
The pipe between the scheduler and the child process had to remain open for the full
lifetime of the child, which could span several SIGSTOP/SIGCONT cycles. Closing the
write end in the parent before the first fork was sufficient; the read end was kept in the
\texttt{Task} struct and only closed after the child exited and all output was drained.

\textbf{Detecting completion at the quantum boundary.}
The polling loop exits when \texttt{elapsed >= quantum}, but if the child process exits
at exactly the same moment, \texttt{waitpid(WNOHANG)} is not called again and the
completion is missed. An additional \texttt{waitpid} call was added immediately after the
loop exits to catch this edge case and avoid an unnecessary extra scheduling round.

\textbf{Preemption notification from client thread to scheduler thread.}
Client threads call \texttt{scheduler\_add\_task()} while the scheduler thread may be
sleeping inside \texttt{run\_program\_slice()}. A simple \texttt{preempt\_flag} field in
the queue, read every 200\,ms by the polling loop, was sufficient. More complex mechanisms
such as signals sent to the scheduler thread were considered but rejected as they would
have complicated signal mask management across threads.

\textbf{Sending results on the correct socket.}
In Phase 3, each client thread was responsible for its own socket. In Phase 4, the
scheduler thread sends results directly to client sockets. This required storing
\texttt{client\_fd} in the \texttt{Task} struct and carefully clearing the pipe file
descriptor in the task slot before releasing the mutex, to avoid a double-close if the
client disconnected concurrently.

\textbf{Avoiding a tight spin when the queue is empty.}
Early versions of the scheduler polled the queue in a busy loop. Replacing the loop with
\texttt{pthread\_cond\_wait()} on \texttt{has\_task} eliminated the CPU waste; client
threads signal the condition variable via \texttt{pthread\_cond\_signal()} in
\texttt{scheduler\_add\_task()}.

% ═════════════════════════════════════════════════════════════════════════════
\section{Division of Tasks}
% ═════════════════════════════════════════════════════════════════════════════

Development responsibilities were divided to match each team member's strengths and focus
areas for this phase.

\textbf{Bernard Gharbin (bg2696)} was primarily responsible for the design and
implementation of \texttt{scheduler.c} and \texttt{scheduler.h}. This included: the
definition of all data structures (\texttt{Task}, \texttt{TaskQueue}, \texttt{HistEntry},
\texttt{TaskState}); the core scheduling loop in \texttt{scheduler\_run()}; the SRJF
task selection algorithm in \texttt{select\_next\_task()} including the no-consecutive
rule and FCFS tie-breaking; the preemption mechanism using \texttt{SIGSTOP} and
\texttt{SIGCONT}; the output pipe management across multiple scheduling slices in
\texttt{fork\_program()} and \texttt{run\_program\_slice()}; the Gantt-chart history
recording and summary printing; and the client-disconnect cleanup logic in
\texttt{scheduler\_remove\_client()}.

\textbf{Bismark Buernortey Buer (bb3621)} was primarily responsible for
\texttt{demo.c}, the modifications to \texttt{server.c} and \texttt{client.c}, and the
updated \texttt{Makefile}. This included: implementing the \texttt{demo.c} simulation
program with correct per-second output and flush behaviour; adding the
``\texttt{Connected to a server}'' message to \texttt{client.c}; adding
\texttt{classify\_command()} to the server to distinguish shell commands from programs
and extract burst times; integrating \texttt{scheduler\_add\_task()} and
\texttt{scheduler\_remove\_client()} into the client thread lifecycle; spawning and
detaching the scheduler thread at server startup; updating the \texttt{Makefile} to add
the \texttt{demo} target and link \texttt{scheduler.o} into the server binary; and
conducting all three required test scenarios on the remote Linux server, capturing output
and screenshots, and authoring this report.

Both team members collaborated on integration testing across multiple terminal windows,
debugging timing issues in the polling loop, and reviewing all scheduling decisions
against the project specification. As in previous phases, the project was developed using
GitHub with a structured branching workflow; each member worked in a separate feature
branch and submitted pull requests for review before merging into the main branch.

% ═════════════════════════════════════════════════════════════════════════════
\section{References}
% ═════════════════════════════════════════════════════════════════════════════

\begin{enumerate}
    \item Linux manual pages for \texttt{pthread\_create}, \texttt{pthread\_mutex\_lock},
          \texttt{pthread\_cond\_wait}, \texttt{pthread\_cond\_signal},
          \texttt{sem\_open}, \texttt{sem\_wait}, \texttt{sem\_post},
          \texttt{fork}, \texttt{pipe}, \texttt{dup2}, \texttt{waitpid},
          \texttt{kill}, \texttt{SIGSTOP}, \texttt{SIGCONT}, \texttt{nanosleep},
          and \texttt{execvp} were referenced throughout development.

    \item The Phase 1, 2, and 3 implementations (\texttt{parse.c}, \texttt{execute.c},
          \texttt{shell.c}, \texttt{server.c}, \texttt{client.c}) served as the
          cumulative foundation for Phase 4, with all modules reused unchanged unless
          otherwise noted.

    \item Course lecture materials and laboratory examples on CPU scheduling algorithms
          --- including Round-Robin, Shortest-Job-First, and preemptive scheduling ---
          were consulted to inform the combined SRJF + RR design and the choice of
          quantum values.
\end{enumerate}

\end{document}
